% !TEX root = ../main.tex

\section{Lesson 02 --- 25/09/2026}

\subsection{Examples on balls}
We start the lesson with some examples on balls.

\begin{example}[Exercise 1.1]
Find the unit ball $B(0,1)$ w.r.t.\ $\norm{\cdot}_1$, $\norm{\cdot}_\infty$ and $\norm{\cdot}_A$, where
\[ A = \begin{bmatrix}
    4 & 0 \\
    0 & 1
\end{bmatrix}.
\]
\begin{itemize}
    \item For $\norm{x}_1 = |x_1|+|x_2|$ the ball is $B(0,1) = \{ x : |x_1|+|x_2|\le1\}$. Since
    \[
     |x_1|+|x_2| =
      \max\{x_1 + x_2,\ x_1 - x_2,\ -x_1 + x_2,\ -x_1 - x_2\},
    \]
    all four expressions must be $\le 1$, so the ball is a square rotated by 45 degrees, with vertices $(\pm 1, 0)$, $(0, \pm 1)$.
    \item For $\norm{x}_\infty = \max\{|x_1|, |x_2|\}$ the ball is $\{x : |x_1| \le 1,\ |x_2| \le 1\}$, i.e.\ the square $[-1,1]^2$.
    \item For $\norm{x}_A = \sqrt{x^TAx}$ we have
    \[
        x^TAx = \begin{bmatrix}x_1 & x_2\end{bmatrix}\begin{bmatrix}
        4 & 0 \\
        0 & 1
    \end{bmatrix}\begin{bmatrix}x_1 \\ x_2\end{bmatrix} = 4x_1^2 + x_2^2 \le 1,
    \]
    which is an ellipse.
\end{itemize}
All of them are convex sets.
\end{example}

The fact that these balls are all convex is not a coincidence.

\begin{example}[Exercise 1.4]
Prove that $B(0,1)=\{x \in \R^n:\norm{x}\le1\}$ is convex for any norm.

Let $x,y \in B(0,1)$; we have to prove that $\alpha x + (1-\alpha)y \in B(0,1)$ for all $\alpha \in[0,1]$. By the triangle inequality and homogeneity,
\[
    \norm{\alpha x + (1-\alpha)y} \le \norm{\alpha x} + \norm{(1-\alpha)y} = \alpha\norm{x}+ (1-\alpha)\norm{y} \le \alpha + (1-\alpha) = 1.
\]
\end{example}

\subsection{Operations that preserve convexity}

\subsubsection{Algebraic operations}

\begin{lemma}
If $C_1$ and $C_2$ are convex, then $C_1 + C_2\coloneqq \{x+y: x \in C_1, y \in C_2\}$ is convex.
If $C$ is convex and $\alpha \in \R$, then $\alpha C \coloneqq \{\alpha x : x \in C\}$ is convex.
Consequently, $C_1 - C_2\coloneqq \{x-y: x \in C_1, y \in C_2\}$ is convex.
\end{lemma}

\subsubsection{Topological operations}

\begin{lemma}
If $C$ is convex then $\mathrm{cl}(C)$ is convex.
If $C$ is convex then $\mathrm{int}(C)$ is convex, provided that $\mathrm{int}(C)\ne \emptyset$.
\end{lemma}

The interior of a convex set can be empty even when the set is ``reasonable'' (e.g.\ a segment in $\R^2$). For this reason it is useful to look at the interior relative to the smallest affine set containing $C$. We denote by $\mathrm{aff}(C)$ the smallest affine set containing $C$.

\begin{definition}[Relative interior]
Let $C \subseteq \R^n$ be a convex set. The \emph{relative interior} of $C$ is
\[
    \mathrm{ri}(C)=\{x \in C : \exists\,\epsilon > 0\ \text{s.t.}\ \mathrm{aff}(C) \cap B(x,\epsilon)\subseteq C\}.
\]
\end{definition}

\begin{example}
Let $C \coloneqq \{(x_1, x_2) \in \R^2 : 1 \le x_1 \le 3,\ x_2 = 0\}$. Then $\mathrm{ri}(C) = \{(x_1, x_2) : 1 < x_1 < 3,\ x_2 = 0\}$, while $\mathrm{int}(C) = \emptyset$. If $C = \{\bar{x}\}$, then $\mathrm{ri}(C) = C$.
\end{example}

\begin{theorem}
If $C$ is a nonempty convex set in $\R^n$, then $\mathrm{ri}(C)$ is a nonempty convex set.
\end{theorem}

\subsubsection{Separation of convex sets}

The relative interior is what is needed to state when two convex sets can be separated by a hyperplane.

Two sets $A, B \subseteq \R^n$ are \emph{linearly separable} if there exist $a \in \R^n$, $a \ne 0$, and $\beta \in \R$ such that
\[
    a^Tx \ge \beta \quad \forall x \in A, \qquad a^Tx \le \beta \quad \forall x \in B.
\]
The separation is \emph{proper} if the strict inequality holds for at least one $x \in A \cup B$.

\begin{theorem}[Separation]
Let $A, B$ be nonempty convex sets in $\R^n$. Then $A$ and $B$ are properly linearly separable if and only if
\[
    \mathrm{ri}(A) \cap \mathrm{ri}(B) = \emptyset.
\]
In particular, two disjoint convex sets are always properly linearly separable.
\end{theorem}

\begin{example}
Let $A \coloneqq \{x \in \R^2 : 1 \le x_1 \le 2,\ x_2 = 0\}$ and $B \coloneqq \{x \in \R^2 : 2 \le x_1 \le 4,\ x_2 = 0\}$. Then $\mathrm{ri}(A) \cap \mathrm{ri}(B) = \emptyset$ and the sets are properly separable by the hyperplane $x_1 = 2$.
\end{example}

\subsubsection{Affine transformations}

\begin{theorem}
Let $f: \R^n \to \R^m$ be affine, i.e.\ $f(x) = Ax + b$, with $A \in \R^{m \times n}$, $b \in \R^m$.
\begin{itemize}
    \item If $C \subseteq \R^n$ is convex, then $f(C) = \{f(x) : x \in C\}$ is convex.
    \item If $C \subseteq \R^m$ is convex, then $f^{-1}(C) = \{x \in \R^n : f(x) \in C\}$ is convex.
\end{itemize}
\end{theorem}

Examples of affine maps are the scaling $f(x) = \alpha x$, the translation $f(x) = x + b$ and the rotation $f(x) = \begin{bmatrix} \cos\theta & -\sin\theta\\ \sin\theta & \cos\theta\end{bmatrix} x$.
